\documentclass[compress, xcolor={svgnames}, aspectratio=169, smaller, t]{beamer}
\newcommand{\CommonPath}{../../Common}
\newcommand{\coursename}{Intro to DL - S2}
\usepackage{\CommonPath/slidespreamble}
\usepackage{\CommonPath/math_commands}
% \setbeameroption{show notes}
\setbeameroption{hide notes}

\title{Introduction to Deep Learning}
\subtitle{Session 2: Perceptron, MLPs \& Activations \\[1em]
\small March 2026}
\author{Magda Gregorová - \href{mailto:magda.gregorova@thws.de}{magda.gregorova@thws.de}}
\institute{
\includegraphics[width=0.3\textwidth]{thws-logo_vert_en_orange-rgb.png}}
\hypersetup{
pdftitle={Introduction to Deep Learning},
pdfauthor={Magda Gregorová},
pdfkeywords={perceptron, MLP, activation functions, backpropagation}
}

\date{{\color{skyblue1}\footnotesize Licensed according to
\href{https://creativecommons.org/licenses/by-sa/4.0}{\color{skyblue1}CC-BY-SA 4.0}}}

\begin{document}

\begin{frame}[plain]
    \titlepage
\end{frame}

\begin{frame}{Lecture Overview}

\tableofcontents

\end{frame}

\section{Supervised Learning}

\begin{frame}{Why Machine Learning}

\vskip 2em

\begin{columns}[c]
  \begin{column}{0.25\textwidth}
\centering
    \structure{Traditional Software}\\[1em]
    \includegraphics[width=4cm]{TraditionalPipeline.excalidraw}
  \end{column}
  \begin{column}{0.4\textwidth}
\centering
\tikz\draw[-Stealth, thick, thwspetrol] (0,0) -- (3.5,0);\\[0.5em]
\scriptsize{
adapt to changing conditions\\
too many variations to capture\\
too complex to understand fully\\
too expensive to compute}
  \end{column}
  \begin{column}{0.25\textwidth}
\centering
    \structure{Machine Learning}\\[1em]
    \includegraphics[width=3.3cm]{MLPipeline.excalidraw}
  \end{column}
  \begin{column}{0.05\textwidth}
\centering
    \
  \end{column}
\end{columns}

\end{frame}

\begin{frame}{Supervised Learning}

\vskip 2em

\structure{Setup:} data pairs $\dataset = \{(\vx^{(i)}, y^{(i)})\}_{i=1}^{\nsamples}$ with $\vx \in \inspace$, $y \in \outspace$

\vskip 0.8em

\structure{Assumption:} exists true relationship $f^*: \inspace \to \outspace, \ y = f^*(\vx)$, \alert{but $f^*$ is never observed!}

\vskip 0.8em

\structure{Goal:} learn $f_\params: \inspace \to \outspace$ from $\dataset$ such that $f_\params \approx f^*$
\alert{$\Rightarrow$ predict on \textbf{unseen} inputs $\vx \notin \dataset$}

\vskip 2em

\centering
% Common/tikz/supervised_learning.tex
\begin{tikzpicture}[scale=0.9]
\tikzstyle{box} = [draw, rounded corners, thick, minimum width=2.8cm, minimum height=1.2cm, align=center, font=\large]
\tikzstyle{lbl} = [font=\small]
\tikzstyle{math} = [color=carnelian, font=\normalsize]
\tikzstyle{annot} = [color=thwsblue, font=\large]

% Three main boxes
\node[box] (ins) at (0, 0) {$x \in \inspace$};
\node[box, fill=thwspetrol!20, draw=thwsblue] (model) at (4.5, 0) {$f_\theta \approx f^*$};
\node[box] (outs) at (9, 0) {$y \in \outspace$};

% Arrows between boxes
\draw[-Stealth, thick] (ins.east) -- (model.west);
\draw[-Stealth, thick] (model.east) -- (outs.west);

% Labels above boxes
\node[annot] [above=0.1 of ins] {Inputs};
\node[annot] [above=0.1 of model] {Model};
\node[annot] [above=0.1 of outs] {Outputs};

% Math labels below boxes
\node[math] (isp) [below=0.3 of ins] {input space};
\node[math] (osp) [below=0.3 of outs] {output space};
\node[math] (param) [below=0.3 of model] {parameters};

% Arrows from math labels up to boxes
\draw[-Stealth, thick, carnelian] (isp) -- (0.3, -0.2);
\draw[-Stealth, thick, carnelian] (9, -1.1) -- (9.3, -0.2);
\draw[-Stealth, thick, carnelian] (4.5, -1.1) -- (4.2, -0.2);

\end{tikzpicture}

\vskip 0.8em

\structure{Notation:} $f_\params(\vx) = f(\vx; \params) = \yhat$

\end{frame}

\begin{frame}{Supervised Learning - Basic Tasks}

\vskip 1em

\begin{columns}
    \begin{column}{0.5\textwidth}
\structure{Classification (binary)}
\begin{itemize}
    \item Output: discrete label $\outspace = \{0, 1\}$
    \item Goal: find $f_\params: \inspace \to \{0, 1\}$
    \item Predict category for input
\end{itemize}
        \vskip 0.5em
        \centering
        % Common/tikz/classification_2d.tex
\begin{tikzpicture}[scale=0.9]
\tikzstyle{posit} = [circle, draw, fill=thwsorange!60, minimum size=0.1cm, inner sep=0pt]
\tikzstyle{neg} = [circle, draw, fill=thwsblue!40, minimum size=0.1cm, inner sep=0pt]
\tikzstyle{annot} = [font=\scriptsize]

% Axes
\draw[-Stealth, thick] (-0.3, 0) -- (3.8, 0) node[right, annot] {$x_1$};
\draw[-Stealth, thick] (0, -0.3) -- (0, 3.2) node[above, annot] {$x_2$};

% Classificaiton boundary
\onslide<2->{
\draw[thick, thwsorange] (0.0, 0.3) -- (3.6, 2.6) node[annot, anchor=south] {separating line};}

% posititive class
\node[posit] at (0.5, 2.4) {};
\node[posit] at (0.8, 2.9) {};
\node[posit] at (0.4, 2.0) {};
\node[posit] at (1.1, 2.6) {};
\node[posit] at (0.7, 3.0) {};

% Negative class
\node[neg] at (2.6, 0.5) {};
\node[neg] at (3.0, 0.9) {};
\node[neg] at (3.2, 0.4) {};
\node[neg] at (2.4, 1.0) {};
\node[neg] at (2.9, 1.3) {};

% Labels
\node[annot, thwsorange] at (1.7, 2.9) {$y=1$};
\node[annot, thwsblue]   at (1.8, 0.5) {$y=0$};


\end{tikzpicture}
    \end{column}
    \begin{column}{0.5\textwidth}
\structure{Regression}
\begin{itemize}
    \item Output: continuous value $\outspace = \R$
    \item Goal: find $f_\params: \inspace \to \R$
    \item Predict numerical value for input
\end{itemize}
        \vskip 0.5em
        \centering
        % Common/tikz/regression_1d.tex
\begin{tikzpicture}[scale=0.9]
\tikzstyle{pt}    = [circle, draw, fill=thwsblue!40, minimum size=0.1cm, inner sep=0pt]
\tikzstyle{annot} = [font=\scriptsize]

% Axes
\draw[-Stealth, thick] (-0.3, 0) -- (3.8, 0) node[right, annot] {$x$};
\draw[-Stealth, thick] (0, -0.3) -- (0, 3.2) node[above, annot] {$y$};

% % Regression line
% \onslide<2->{
% \draw[thick, thwsorange] (0.1, 0.4) -- (3.6, 2.6) node[annot, right, anchor=south] {regression line};}

% Data points scattered around curve
\node[pt] at (0.3, 0.6) {};
\node[pt] at (0.7, 0.9) {};
\node[pt] at (1.0, 0.8) {};
\node[pt] at (1.3, 1.2) {};
\node[pt] at (1.6, 1.5) {};
\node[pt] at (1.9, 1.4) {};
\node[pt] (yx) at (2.2, 1.8) {};
\node[pt] at (2.5, 2.2) {};
\node[pt] at (2.8, 2.1) {};
\node[pt] at (3.2, 2.2) {};
\node[pt] at (3.5, 2.5) {};

% y point
\draw[dotted, thick] (yx) -- ++(0, -1.9) node[below, annot] {$x^{(i)}$};
\draw[dotted, thick] (yx) -- ++(-2.2, 0) node[left, annot] {$y^{(i)}$};;

% Label
% \node[annot, thwsorange] at (2.8, 2.8) {$f(x)$};
\end{tikzpicture}
    \end{column}
\end{columns}
\end{frame}

\begin{frame}{Regression}
\begin{columns}
    \begin{column}{0.55\textwidth}
        \begin{itemize}
            \item Output is continuous: $y \in \R$
            \item Given data $\dataset = \{(x^{(i)}, y^{(i)})\}_{i=1}^{\nsamples}$
            \item Goal: find $f(x; \params)$ that fits the data and generalizes
            \vskip 1em
            \item Simplest model: a line
            \[ f(x; \params) = wx + b, \quad \params = (w, b) \]
        \end{itemize}
    \end{column}
    \begin{column}{0.42\textwidth}
        \centering
        % Common/tikz/regression_1d.tex
\begin{tikzpicture}[scale=0.9]
\tikzstyle{pt}    = [circle, draw, fill=thwsblue!40, minimum size=0.1cm, inner sep=0pt]
\tikzstyle{annot} = [font=\scriptsize]

% Axes
\draw[-Stealth, thick] (-0.3, 0) -- (3.8, 0) node[right, annot] {$x$};
\draw[-Stealth, thick] (0, -0.3) -- (0, 3.2) node[above, annot] {$y$};

% % Regression line
% \onslide<2->{
% \draw[thick, thwsorange] (0.1, 0.4) -- (3.6, 2.6) node[annot, right, anchor=south] {regression line};}

% Data points scattered around curve
\node[pt] at (0.3, 0.6) {};
\node[pt] at (0.7, 0.9) {};
\node[pt] at (1.0, 0.8) {};
\node[pt] at (1.3, 1.2) {};
\node[pt] at (1.6, 1.5) {};
\node[pt] at (1.9, 1.4) {};
\node[pt] (yx) at (2.2, 1.8) {};
\node[pt] at (2.5, 2.2) {};
\node[pt] at (2.8, 2.1) {};
\node[pt] at (3.2, 2.2) {};
\node[pt] at (3.5, 2.5) {};

% y point
\draw[dotted, thick] (yx) -- ++(0, -1.9) node[below, annot] {$x^{(i)}$};
\draw[dotted, thick] (yx) -- ++(-2.2, 0) node[left, annot] {$y^{(i)}$};;

% Label
% \node[annot, thwsorange] at (2.8, 2.8) {$f(x)$};
\end{tikzpicture}
    \end{column}
\end{columns}
\end{frame}

\begin{frame}{The linear model}
\begin{columns}
    \begin{column}{0.55\textwidth}
        \begin{itemize}
            \item $f(x; \params) = wx + b$, \quad $\params = (w, b)$
            \vskip 0.5em
            \item<2-> Works well when data is linear
            \vskip 0.5em
            \item<3-> Fails when data has curvature
            \vskip 0.5em
            \item<3-> We need a more expressive $f(\cdot\,; \params)$
        \end{itemize}
    \end{column}
    \begin{column}{0.42\textwidth}
        \centering
        \texttt{[linear\_fit.tex]}
    \end{column}
\end{columns}
\end{frame}

\begin{frame}{Piecewise linear functions}
\begin{columns}
    \begin{column}{0.55\textwidth}
        \begin{itemize}
            \item What if we allow the function to bend?
            \vskip 0.5em
            \item<2-> A piecewise linear function:
            \begin{itemize}
                \item linear between breakpoints
                \item can approximate any continuous function
            \end{itemize}
            \vskip 0.5em
            \item<3-> More breakpoints = more expressive $f(\cdot\,; \params)$
        \end{itemize}
    \end{column}
    \begin{column}{0.42\textwidth}
        \centering
        \texttt{[piecewise\_linear.tex]}
    \end{column}
\end{columns}
\end{frame}

\begin{frame}{ReLU as a building block}
\begin{columns}
    \begin{column}{0.55\textwidth}
        \begin{itemize}
            \item Rectified Linear Unit:
            \[ \text{ReLU}(x) = \max(0, x) \]
            \item<2-> A single ReLU introduces one bend
            \vskip 0.5em
            \item<3-> Linear combination of ReLUs:
            \[ f(x; \params) = \sum_j v_j \, \text{ReLU}(w_j x + b_j) \]
            \item<3-> Each term contributes one breakpoint
            \vskip 0.5em
            \item<4-> More neurons = more bends = more expressive
        \end{itemize}
    \end{column}
    \begin{column}{0.42\textwidth}
        \centering
        \texttt{[relu\_combination.tex]}
    \end{column}
\end{columns}
\end{frame}

\begin{frame}{This is a neural network}
\begin{columns}
    \begin{column}{0.55\textwidth}
        \begin{itemize}
            \item Each hidden neuron computes:
            \[ z_j^{(1)} = w_j x + b_j, \quad a_j^{(1)} = \text{ReLU}(z_j^{(1)}) \]
            \item Output layer combines:
            \[ f(x; \params) = \sum_j v_j a_j^{(1)} + c \]
            \vskip 0.5em
            \item<2-> In matrix form:
            \begin{align*}
                \vz^{(1)} &= \mW^{(1)}\vx + \vb^{(1)} \\
                \va^{(1)} &= \text{ReLU}(\vz^{(1)}) \\
                f(\vx; \params) &= \mW^{(2)}\va^{(1)} + \vb^{(2)}
            \end{align*}
            \item<2-> Parameters: $\params = \{\mW^{(1)}, \vb^{(1)}, \mW^{(2)}, \vb^{(2)}\}$
        \end{itemize}
    \end{column}
    \begin{column}{0.42\textwidth}
        \centering
        \texttt{[mlp\_1hidden.tex]}
    \end{column}
\end{columns}
\end{frame}

\begin{frame}{Universal approximation}
    \begin{block}{Universal Approximation Theorem}
        Let $\act: \R \to \R$ be a non-constant, bounded, continuous activation function.
        Then for any continuous function $g: [0,1]^d \to \R$ and any $\varepsilon > 0$,
        there exists a 1-hidden-layer neural network $f(\cdot\,; \params)$ such that:
        \[ \sup_{\vx \in [0,1]^d} |f(\vx; \params) - g(\vx)| < \varepsilon \]
    \end{block}
    \vskip 1em
    \begin{columns}
        \begin{column}{0.48\textwidth}
            \structure{What it tells us}
            \begin{itemize}
                \item MLPs are universal approximators
                \item Enough neurons can fit any continuous function
            \end{itemize}
        \end{column}
        \begin{column}{0.48\textwidth}
            \structure{What it doesn't tell us}
            \begin{itemize}
                \item How many neurons are needed
                \item How to find the parameters $\params$
                \item Whether the model will generalize
            \end{itemize}
        \end{column}
    \end{columns}
\end{frame}

\begin{frame}{Classification}
\begin{columns}
    \begin{column}{0.55\textwidth}
        \begin{itemize}
            \item Output is discrete: $y \in \{0, 1\}$
            \vskip 0.5em
            \item Linear classifier: decision boundary is a hyperplane
            \[ \hat{y} = \mathbb{1}[f(\vx; \params) \geq 0] \]
            \vskip 0.5em
            \item<2-> Not all problems are linearly separable
            \begin{itemize}
                \item<2-> XOR: four points, no linear boundary works
            \end{itemize}
            \vskip 0.5em
            \item<3-> Neural network with ReLU: non-linear decision boundary
            \item<3-> Same architecture as for regression
        \end{itemize}
    \end{column}
    \begin{column}{0.42\textwidth}
        \centering
        \only<1>{\texttt{[linear\_boundary.tex]}}
        \only<2->{\texttt{[xor\_boundary.tex]}}
    \end{column}
\end{columns}
\end{frame}


\begin{frame}{Activation functions}
\begin{columns}
    \begin{column}{0.48\textwidth}
        \structure{What makes a good activation?}
        \begin{itemize}
            \item Non-linear — otherwise depth buys nothing
            \item Differentiable — needed for backpropagation
            \item Computationally cheap
        \end{itemize}
        \vskip 1em
        \structure{Common choices}
        \begin{itemize}
            \item \textbf{Sigmoid}: $\sigma(z) = \frac{1}{1+e^{-z}} \in (0,1)$
            \item \textbf{Tanh}: $\tanh(z) \in (-1, 1)$
            \item \textbf{ReLU}: $\max(0, z)$ — sparse, fast, but dying neurons
        \end{itemize}
        \vskip 0.5em
        \structure{Dying ReLU problem}
        \begin{itemize}
            \item Neurons with $z < 0$ output 0 and have zero gradient
            \item Can become permanently inactive during training
        \end{itemize}
    \end{column}
    \begin{column}{0.48\textwidth}
        \centering
        \texttt{[activation\_functions.tex]}
    \end{column}
\end{columns}
\end{frame}

\begin{frame}{Notation summary}
\begin{center}
\renewcommand{\arraystretch}{1.4}
\begin{tabular}{cll}
    \hline
    \textbf{Symbol} & \textbf{Meaning} & \textbf{Dimensions} \\
    \hline
    $\vx$ & input & $\R^{d}$ \\
    $y$ & target output & $\R$ or $\{0,1\}$ \\
    $\params$ & all parameters & — \\
    $\mW^{(l)}$ & weight matrix, layer $l$ & $\R^{d_{l} \times d_{l-1}}$ \\
    $\vb^{(l)}$ & bias vector, layer $l$ & $\R^{d_{l}}$ \\
    $\vz^{(l)}$ & pre-activation, layer $l$ & $\R^{d_{l}}$ \\
    $\va^{(l)}$ & post-activation, layer $l$ & $\R^{d_{l}}$ \\
    $\act$ & activation function & $\R \to \R$ \\
    $L$ & number of hidden layers & — \\
    \hline
\end{tabular}
\end{center}
\vskip 0.5em
\begin{center}
    \structure{Forward pass:} \quad
    $\vz^{(l)} = \mW^{(l)}\va^{(l-1)} + \vb^{(l)}$ \quad
    $\va^{(l)} = \act(\vz^{(l)})$ \quad
    $\va^{(0)} = \vx$
\end{center}
\end{frame}

\end{document}